\section{Responsibility} 
It is a well known fact that most programmers are poor technical writers. That's one of the reasons why programmers rarely write the public documentation for their products. According to a rumor, the other reason is that programmers are too expensive to write simple documentation. The fact is, however, that skilled technical writers can easily command salaries on par with those of most top programmers.

Nevertheless, there is no excuse why programmers do not write proper documentation comments into their source code. They are the only people that could write these comments because technical writers usually do neither have the time nor the required skills to analyse the code and to extract the information from it that is necessary for the documentation – not to mention that the technical writers often do not have (write) access to the source code and are therefore being technically not able to add the documentation comments.

Furthermore there are two main target groups for the documentation that is generated from the documentation comments: the first group are the peer developers writing code interoperating with the already existing one, and the second group are the maintenance programmers that have to fix what's broken.

A third group are the testers that are designing their tests based on the generated documentation, but they also use the design documents and the public documentation as their sources.

This alone suggests already that documentation comments should be written by the programmers themselves – programmers writing for programmers!

If your product is a \nameref{sec:FunctionLibrary}, the generated documentation may be all that you provide to your clients; same for a \nameref{sec:FeatureLibrary}, although that usually requires some documentation \textit{supplementing} the generated documentation, not substituting it.

However, even when the generated documentation is intended solely for internal use, it should be created with the same care as documentation for an external audience; in fact, internal documentation should be even more comprehensive and detailed, capturing elements from the deepest recesses of the \bdq{engine room} – things an external customer is better off never seeing. Another reason why the documentation comments belong to the responsibilities of the programmers.

%==============================================================================
% End of file!!
%==============================================================================
